\chapter{Plano de trabalho consolidado}

\section{Identificacao}

\begin{longtable}{@{}p{0.31\linewidth}p{0.61\linewidth}@{}}
\toprule
\textbf{Campo} & \textbf{Valor} \\
\midrule
Titulo & Estudo de eficiencia computacional de operacoes de Transformers sob diferentes configuracoes numericas e estruturais \\
Estudante & Lucas Cerqueira Galvao \\
Curso & Ciencia da Computacao \\
Grupo de pesquisa & LPCIAD - Laboratorio de Projetos de Circuitos Integrados Analogicos e Digitais \\
Orientacao & Prof. Me. Walter Silva Oliveira \\
Vigencia original & Maio a setembro de 2026 \\
Dedicacao registrada & 8 horas por semana, 4 meses, 80 horas-aula de atividade curricular de pesquisa \\
Projeto vinculado & Registro Ipeci 20240004 - Projeto e Otimizacao de Circuitos Integrados utilizando Inteligencia Artificial \\
\bottomrule
\end{longtable}

O formulario institucional submetido permanece preservado em
\sourcepath{docs/plano_trabalho/PlanoTrabalhoEP_LucasGalvao.docx}. Esta secao e
a versao academica editavel em LaTeX e omite instrucoes de preenchimento,
assinaturas e dados pessoais que nao interferem no metodo cientifico.

\section{Resumo}

A arquitetura Transformer consolidou-se como alicerce de sistemas contemporaneos
de inteligencia artificial, mas impoe custo significativo de execucao, memoria e
operacoes aritmeticas durante a inferencia. O estudo analisa como configuracoes
numericas e estruturais afetam operacoes de autoatencao e projecoes lineares
densas. Um baseline sem otimizacao e comparado com pruning por magnitude e
quantizacao linear INT8. A avaliacao combina latencia, pico de memoria e FLOPs
analiticos com MSE, MAE, coeficiente de determinacao e similaridade de cosseno,
buscando caracterizar o compromisso entre reducao de custo e preservacao da
saida.

\section{Objetivo geral}

Analisar como diferentes configuracoes numericas e estruturais afetam a
eficiencia computacional de operacoes centrais da arquitetura Transformer em
cenarios controlados de inferencia.

\section{Objetivos especificos}

\begin{enumerate}
  \item implementar, em Python, um modelo computacional simplificado com foco em autoatencao e projecao linear densa;
  \item construir cenarios comparaveis de baseline, pruning por magnitude e quantizacao linear;
  \item mensurar latencia, memoria e estimativas analiticas de operacoes;
  \item avaliar as saidas por MSE, MAE, $R^2$ e similaridade de cosseno;
  \item interpretar os compromissos entre reducao de complexidade e preservacao da acuracia numerica.
\end{enumerate}

\section{Contrato experimental}

O protocolo principal usa Python 3.11, PyTorch/CUDA, NumPy, Pandas,
scikit-learn e Matplotlib. A grade formal tem lote 1, comprimentos de sequencia
$L \in \{64,128,256\}$ e dimensoes $D \in \{128,256,512\}$. Cada combinacao
mede baseline, pruning de 50\% e quantizacao INT8, com 20 iteracoes de
aquecimento, 50 medicoes validas e duas execucoes independentes. A semente
principal e 42; entradas seguem $\mathcal{N}(0,1)$, pesos usam inicializacao
Xavier normal ajustada a dimensao e o bias e zero.

\section{Infraestrutura}

A coleta principal foi planejada para uma NVIDIA GeForce RTX 4070 Ti Super de
16 GB. TensorFlow/Keras participa apenas do portao de equivalencia numerica em
CPU; o benchmark de desempenho permanece no caminho PyTorch/CUDA. A pesquisa
nao requer laboratorio especifico nem material de consumo, alem da infraestrutura
computacional ja disponivel.
